\paragraph{到达间隔包含钟慢与传播路程两部分}
在观察者惯性系中，光源沿视线以 $v>0$ 远离。两次发射在光源自身钟上相隔 $\Delta\tau=1/f_0$，在观察者系中相隔 $\Delta t=\gamma\Delta\tau$。第二次发射时光源又远离了 $v\Delta t$，第二束光还要多走这段路，所以
\[
 \Delta t_{\rm rec}=\Delta t+\frac{v\Delta t}{c}
 =\gamma(1+\beta)\Delta\tau,\qquad
 r=\frac{f_0}{f_{\rm rec}}=\gamma(1+\beta).
\]
平方并用 $\gamma^2=(1-\beta^2)^{-1}$，即可得到前面的 $r^2=(1+\beta)/(1-\beta)$。仅使用时间延缓会漏掉传播路程变化，也就无法分辨靠近与远离。
对沿视线靠近的源，可在同一约定下取 $\beta<0$，此时 $r<1$，反解式自然给出负速度；$r=1$ 给出零径向速度，$r\to\infty$ 时 $\beta\to1^-$，不会反推出超光速运动。非沿视线运动还须使用含观测角的完整多普勒关系。
